\chapter{Aan de slag met BASIC}
\label{cha:basic-getting-started}

Op de MEGA65 is het mogelijk om in veel talen te programmeren, maar de meeste mensen beginnen met BASIC. Dat is begrijpelijk, want BASIC is de afkorting van Beginner's All-purpose Symbolic Instruction Code: de taal is gemaakt om mensen zoals u de eerste stappen in de wereld van het programmeren te laten zetten!

Nog enkele korte woorden voordat we van start gaan: BASIC is een programmeertaal en heeft, net als een gesproken taal, conventies, grammatica en woordenschat. Gelukkig is BASIC veel sneller en gemakkelijker te leren dan onze ingewikkelde menselijke talen. Als u echter goed oplet, zult u een aantal van die structuren herkennen en dat kan u op weg helpen in de wereld van het programmeren.

Indien u \bookvref{cha:getting-started} nog niet gelezen heeft, is het verstandig dat alsnog te doen, waardoor u zelfverzekerder met de MEGA65-computer kunt omgaan.

Het is ook goed om te weten dat u, als u de MEGA65 echt in de war heeft gebracht, altijd terug kunt keren naar de READY.-prompt: druk gewoon op de resetknop aan de linkerkant van het toetsenbord, of, als dat niet helpt, schakel de computer met de aan/uit-schakelaar aan de linkerkant van het toetsenbord uit en weer in. U hoeft zich geen zorgen te maken over het netjes afsluiten van de computer en soortgelijke idioterie. Het enige dat u moet onthouden, is dat niet-opgeslagen werk verloren gaat wanneer u de computer uitschakelt en weer inschakelt of op de resetknop drukt.

Tot slot: als u niet alle beschrijvingen en toelichtingen bij een voorbeeld begrijpt, geen zorgen! Wij hebben zoveel mogelijk informatie opgenomen, zodat die beschikbaar is wanneer u vragen heeft, problemen tegenkomt of gewoon meer wilt ontdekken. Blader gerust door naar de voorbeelden en probeer zelf van alles uit; u kunt later terugkeren en het nog eens nalezen wanneer u iets te weten wilt komen of een probleem wilt oplossen. En als u geen antwoord op uw probleem vindt, stuur ons dan een bericht! Er zijn ondersteuningsfora voor de MEGA65 op \url{https://mega65.net} en u kunt problemen met deze handleiding melden op:

\url{https://github.com/mega65/mega65-user-guide}

Wij hopen dat u net zoveel plezier beleeft aan het leren programmeren van de MEGA65 als wij hebben gehad bij het maken ervan!

\section{Uw eerste BASIC-programma's}

De MEGA65 is ontworpen om te programmeren! Wanneer u hem inschakelt, heeft de computer een paar seconden nodig om zijn zaken op orde te brengen en daarna toont hij u direct een ``READY.''-prompt en een knipperend blokje dat de cursor wordt genoemd. Zolang de cursor knippert, vertelt hij u dat de computer wacht op invoer. De ``READY.''-melding vertelt u dat de programmeertaal BASIC draait en gereed is om te beginnen met programmeren. U hoeft niet eens een programma te laden -- u kunt gewoon van start gaan.

\needspace{4cm} % Dont allow following paragraph to separate from
                % following element
Typ het volgende in en kijk wat er gebeurt:

\begin{screencode}
HALLO COMPUTER
\end{screencode}

\needspace{4cm} % Dont allow following paragraph to separate from
                % following screenshot

Typ daarvoor gewoon de letters zoals u ze hierboven ziet. De computer staat al in de hoofdlettermodus, dus u hoeft \specialkey{SHIFT} of \specialkey{CAPS\\LOCK} niet ingedrukt te houden. Druk op \specialkey{RETURN} nadat u "HALLO COMPUTER" heeft getypt. De computer geeft u daarmee te kennen dat hij de door u getypte invoerregel moet aanvaarden. Indien u dat doet, ziet u ongeveer de volgende melding:

\screenshotwrap{dutch/images/getting-started/syntax-error.pdf}

  Als u een melding zag die ongeveer lijkt op \screentextwide{SYNTAX ERROR}, dan gefeliciteerd: u bent erin geslaagd om met de computer te communiceren!\index{Fouten!Syntax}\index{SYNTAX ERROR}  Foutmeldingen klinken veel erger dan ze zijn. De MEGA65 gebruikt ze, in het bijzonder de syntaxisfout, om u te laten weten wanneer hij moeite heeft om te begrijpen wat u heeft getypt of wat u in een programma heeft gezet. U hoeft er niet bang voor te zijn; ervaren programmeurs krijgen er voortdurend mee te maken.

  In dit geval raakte de computer in de war, omdat hij het woord ``hallo'' of het woord ``computer'' niet begrijpt. Dat wil zeggen: hij wist niet wat u wilde dat hij deed. In dat opzicht zijn computers tamelijk dom: ze kennen slechts enkele woorden en zijn niet bijzonder vindingrijk in de manier waarop ze die interpreteren.

\needspace{4cm} % Dont allow following paragraph to separate from
                % following screenshot

Laten we het dus nog eens proberen op een manier die de computer wel begrijpt. Typ het volgende in. Dat kunt u direct doen; het maakt niet uit dat de syntaxisfoutmelding nog op het scherm staat. Tegen de tijd dat hij u weer \screentextwide{READY.} toonde, was de computer die alweer vergeten.

\begin{screencode}
PRINT "HALLO COMPUTER"
\end{screencode}

Zorg er opnieuw voor dat u tijdens het intypen geen shift of shift-lock gebruikt. De symbolen rondom het woord \screentextwide{HALLO COMPUTER} zijn dubbele aanhalingstekens. Als u het toetsenbord van een PC gewend bent, zult u wellicht merken dat de toets voor het dubbele aanhalingsteken op een heel andere plaats zit dan u gewend bent: dubbele aanhalingstekens typt u op de MEGA65 door \specialkey{SHIFT} ingedrukt te houden en vervolgens \megakey{2} in te drukken. Vergeet niet op \specialkey{RETURN} te drukken wanneer u klaar bent, zodat de computer weet dat u wilt dat hij iets met uw invoer doet.

Als u tijdens het typen een fout maakt, kunt u \specialkey{INST\\DEL} gebruiken om de fout uit te wissen en te herstellen. U kunt ook de cursor-toetsen gebruiken om op de regel heen en weer te gaan terwijl u de regel bewerkt die u aan het typen bent, maar er schuilt een trucje in wanneer u al een dubbel aanhalingsteken getypt heeft: als u de cursor-toetsen probeert te gebruiken, wordt er een grappig omgekeerd teken weergegeven in plaats van dat de cursor beweegt. Dat komt doordat de computer denkt dat u de verplaatsing van de cursor in de tekst zelf wilt vastleggen, wat erg nuttig en leuk kan zijn en waarover u meer kunt lezen in \bookvref{cha:getting-started}. Voor nu geldt echter: maakt u een fout, druk dan gewoon op \specialkey{RETURN} en typ de verprutste regel opnieuw.

\needspace{4cm} % Dont allow following paragraph to separate from
                % following screenshot
Hopelijk ziet u nu zoiets als dit:

\screenshotwrap{dutch/images/getting-started/print-hello-computer.pdf}

Deze keer zou er geen nieuwe \screentextwide{SYNTAX ERROR}-melding moeten verschijnen. Mocht er echter een of andere foutmelding zijn verschenen, probeer het commando dan gewoon opnieuw in te typen, nadat u goed gekeken heeft waar de fout zou kunnen zitten.

In plaats van een fout zouden we \screentextwide{HALLO COMPUTER} herhaald moeten zien onder de regel die u getypt heeft. De reden dat dit gebeurt, is dat de computer het woord \screentextwide{PRINT} wel degelijk begrijpt. Hij weet dat alles wat er achter het woord \screentextwide{PRINT} staat op het scherm weergegeven moet worden. We moesten \screentextwide{HALLO  COMPUTER} tussen dubbele aanhalingstekens plaatsen om de computer te vertellen dat we willen dat het letterlijk wordt weergegeven.

Hadden we de dubbele aanhalingstekens niet geplaatst, dan zou de computer hebben gedacht dat \screentextwide{HALLO COMPUTER} de naam was van een opgeslagen stukje gegevens. Aangezien we op zo'n plaats geen enkel stukje gegevens hebben opgeslagen, staat er voor de computer een nul en zal de computer dus het getal nul weergeven. Indien de computer een nul of een ander getal weergeeft terwijl u een of andere tekst verwachtte, kan dit de oorzaak zijn.

\needspace{4cm} % Dont allow following paragraph to separate from
                % following screenshot
U kunt het desgewenst zelf proberen en u zou dan zoiets als het volgende moeten zien:

  \screenshotwrap{dutch/images/getting-started/print-hello-computer-no-quotes.pdf}

In de bovenstaande voorbeelden typten we commando's rechtstreeks in en de computer voerde ze onmiddellijk uit nadat u op \specialkey{RETURN} drukte. Het op deze manier typen van commando's wordt derhalve vaak {\em directe modus} genoemd (soms ook wel ``immediate mode'' genoemd).

We kunnen de computer echter ook een lijst van commando's laten onthouden die één voor één uitgevoerd moeten worden. Dat doen we met de nogal fantasieloos benoemde {\em niet-directe modus}. Om de niet-directe modus te gebruiken, zetten we gewoon een getal tussen 0 en 63999 aan het begin van het commando. De computer onthoudt dat commando vervolgens. Anders dan bij een commando in de directe modus, geeft de computer niet opnieuw \screentextwide{READY.} weer. In plaats daarvan verschijnt de cursor gewoon weer op de volgende regel, klaar om meer commando's in te typen.

\needspace{4cm} % Dont allow following paragraph to separate from
                % following screenshot
Laten we dat eens proberen met een eenvoudig programmaatje. Typ de volgende drie invoerregels in:

\begin{screencode}
1 FOR I = 1 TO 10 STEP 1
2 PRINT I
3 NEXT I
\end{screencode}
\index{FOR}
\index{BASIC 65-commando's!FOR}

\needspace{4cm} % Dont allow following paragraph to separate from
                % following screenshot
Indien u dit gedaan heeft, zou het scherm er ongeveer zo uit moeten zien:

\screenshotwrap{images/getting-started/first-steps-for-loop-programme-1.png}

Als dat niet zo is, kunt u het opnieuw proberen. Vergeet niet dat u, wanneer u merkt dat de computer helemaal in de war raakt, gewoon de resetknop kunt indrukken of de aan/uit-schakelaar aan de linkerkant van de computer uit en weer aan kunt zetten om de computer opnieuw op te starten. Dat duurt slechts een paar seconden en doet de MEGA65 geen enkel kwaad.

We hebben de computer gevraagd drie commando's te onthouden, namelijk \screentextwide{FOR I = 1 TO 10 STEP 1},
\screentextwide{PRINT I}
en \screentextwide{NEXT I}. We hebben de computer ook gevraagd in welke volgorde we ze willen uitvoeren: de computer begint met het commando met het laagste nummer en voert vervolgens elk commando met het eerstvolgende hogere nummer uit, totdat hij het einde van de lijst bereikt. Het is dus een beetje als een herinneringslijst voor de computer. Dit is wat we een programma noemen, een beetje zoals het programma van een concert of een theatervoorstelling: het vertelt ons wat er gaat komen en in welke volgorde. Laten we de computer dus vragen dit programma uit te voeren.

Laten we eerst proberen te raden wat er zal gebeuren. Laten we beginnen met het middelste commando, \screentextwide{PRINT I}. We hebben het \screentextwide{PRINT}-commando al gezien en we weten dat het de computer vertelt dat hij dingen op het scherm moet weergeven. Het ding dat hij zal proberen weer te geven is \screentextwide{I}. Net als eerder, omdat er geen dubbele aanhalingstekens rond \screentextwide{I} staan, zal hij proberen een opgeslagen stukje gegevens weer te geven. Het stukje gegevens dat hij probeert weer te geven, is het stukje dat hoort bij \screentextwide{I}.

Indien we een stukje gegevens als dit een naam geven, noemen we dat een {\em variabele}\index{variabele}. Ze worden variabelen genoemd omdat ze kunnen variëren. Dat wil zeggen dat we het stukje gegevens dat hoort bij de variabele genaamd I kunnen vervangen door een ander stukje gegevens. Het oude stukje wordt daardoor vergeten. Als we dus een commando als \screentextwide{LET I = 3} zouden geven, zou dit alles wat er in de variabele \screentextwide{I} was opgeslagen vervangen door het getal 3.

Laten we teruggaan naar ons programma: we weten inmiddels dat het 2\textsuperscript{e} commando zal proberen het stukje gegevens weer te geven dat in de variabele \screentextwide{I} is opgeslagen. Laten we dus naar het eerste commando kijken: \screentextwide{FOR I = 1 TO 10 STEP 1}. Ook al hebben we het \screentextwide{FOR}-commando niet eerder gezien, we kunnen wel een beetje raden hoe het werkt. Het lijkt erop dat het iets in de variabele \screentextwide{I} zal zetten. Dat iets lijkt te maken te hebben met het bereik van de getallen 1 tot en met 10 en met een stap of interval van 1. Wat denkt u dat het zal doen?

Als u geraden heeft dat het de waarden 1, 2, 3, 4, 5, 6, 7, 8, 9 en daarna 10 in de variabele \screentextwide{I} zet, dan verdient u een schouderklopje, want dat is precies wat het doet. Het helpt ons ook het 3\textsuperscript{e} commando te begrijpen, \screentextwide{NEXT I}: dat commando zegt tegen de computer dat hij de volgende waarde in de variabele \screentextwide{I} moet zetten. En hier volgt een beetje magie: wanneer de computer dat doet, gaat hij terug in de lijst met commando's en vervolgt hij vanaf het commando na het \screentextwide{FOR}-commando.

Laten we dit eens samenvoegen: wanneer de computer het eerste commando uitvoert, ontdekt hij dat hij tien verschillende waarden in de variabele \screentextwide{I} moet zetten. Hij begint met het eerste getal erin te zetten, in dit geval het getal 1. De computer gaat vervolgens verder met het tweede commando, dat hem vertelt het stukje gegevens weer te geven dat op dat moment in de variabele \screentextwide{I} is opgeslagen. Dat zal het getal 1 zijn, aangezien dat het laatste was dat de computer erin moest zetten. Daarna gaat de computer door naar het derde commando, dat hem vertelt dat het tijd is om de volgende waarde in de variabele \screentextwide{I} te zetten. De computer gooit dus het getal 1 weg dat zich op dat moment in de variabele \screentextwide{I} bevindt en zet het getal 2 erin, aangezien dat het volgende getal in de lijst is. Vervolgens gaat hij verder vanaf het 2\textsuperscript{e} commando, waardoor de computer de inhoud van de variabele \screentextwide{I} opnieuw weergeeft. Alleen is deze keer het getal 2 het laatst in \screentextwide{I} opgeslagen, dus geeft de computer het getal 2 weer. Dit proces herhaalt zich totdat de computer alle tien waarden heeft weergegeven die het \screentextwide{FOR}-commando aangaf.

\needspace{4cm} % Dont allow following paragraph to separate from
                % following screenshot
Om dit in actie te zien, moeten we de computer het programma van commando's laten uitvoeren dat we hebben ingetypt. We doen dat met het \screentextwide{RUN}-commando. Omdat we willen dat het programma onmiddellijk wordt uitgevoerd, gebruiken we de directe modus. Typ dus gewoon het woord \screentextwide{RUN} en druk op \specialkey{RETURN}. U zou dan een beeld moeten zien dat er ongeveer als volgt uitziet:

\screenshotwrap{images/getting-started/first-steps-for-loop-programme-1-running.png}

U zult hier misschien een paar dingen opmerken:

Ten eerste: de computer vertelt ons weer dat hij \screentextwide{READY.} is, zodra hij klaar is met het uitvoeren van het programma. Dit maakt het voor ons alleen maar makkelijker om te weten wanneer we weer commando's aan de computer kunnen geven.

Ten tweede: toen de computer de onderkant van het scherm bereikte, schoof hij het beeld automatisch omhoog om ruimte te maken. Dat is heel normaal. Wat belangrijk is om te onthouden, is dat de computer alles vergeet wat boven uit het scherm schuift. De enige uitzondering is wanneer u de computer een commando heeft laten onthouden door er een nummer voor te zetten. Ons programma is dus voorlopig veilig. Dat kunnen we zien door het \screentextwide{RUN}-commando nog een paar keer te typen: de programmalijst zal dan boven uit het scherm geschoven zijn, maar we kunnen het programma nog steeds uitvoeren, omdat de computer het heeft onthouden. Probeer het maar! Heeft het gewerkt?

\needspace{4cm} % Dont allow following paragraph to separate from
                % following screenshot
Wilt u het programma met onthouden commando's zien, dan kunt u het \screentextwide{LIST}\index{LIST}\index{BASIC 65-commando's!LIST}-commando gebruiken. Dit commando laat de computer het onthouden programma op het scherm weergeven, zoals in het voorbeeld hier. Wilt u een van de commando's in het programma vervangen, dan kunt u een nieuwe regel typen met hetzelfde nummer als de regel die u wilt wijzigen.

\screenshotwrap{images/getting-started/first-steps-for-loop-programme-1-listing.png}

\needspace{4cm} % Dont allow following paragraph to separate from
                % following screenshot
Om de resultaten bijvoorbeeld allemaal op één regel af te drukken, kunnen we de tweede regel van het programma wijzigen in \screentextwide{PRINT I;} door de volgende regel in te typen en op \specialkey{RETURN} te drukken:



\begin{screencode}
2 PRINT I;
\end{screencode}

\index{PRINT}
\index{BASIC 65-commando's!PRINT}
%\end{tcolorbox}

\needspace{4cm} % Dont allow following paragraph to separate from
                % following screenshot

U kunt controleren of de wijziging is onthouden door het \screentextwide{LIST}-commando opnieuw uit te voeren, zoals we hier zien. Vervolgens kunt u het \screentextwide{RUN}-commando gebruiken om het gewijzigde programma uit te voeren, als volgt:

\screenshotwrap{images/getting-started/first-steps-for-loop-programme-1-modified.png}

Het is heel eenvoudig om uw programma's op deze manier te wijzigen. Naarmate u meer vertrouwd raakt met de werkwijze, zijn er nog twee handige trucs:

Ten eerste kunt u het {\bf LIST}-commando het nummer van een commando geven, of van een regel, zoals commando's ook wel worden aangeduid en het zal alleen die regel van het programma weergeven. U kunt ook een bereik opgeven, gescheiden door een minteken, om alleen een deel van het programma weer te geven, bijv. \screentextwide{LIST 1 - 2} om de eerste twee regels van ons programma te tonen.

Ten tweede kunt u de cursor-toetsen gebruiken om de cursor naar een regel te verplaatsen die al is onthouden en op het scherm wordt weergegeven. Wijzigt u wat u op het scherm ziet en drukt u vervolgens op \specialkey{RETURN} terwijl de cursor op die regel staat, dan leest de BASIC-interpreter de gewijzigde regel in en vervangt hij de oude versie ervan. Het is belangrijk om op te merken dat u, wanneer u meerdere regels van het programma tegelijk wijzigt, op elke gewijzigde regel \specialkey{RETURN} moet drukken. Het is goed gebruik om te controleren of het programma correct is gewijzigd. Gebruik hiervoor opnieuw het {\bf LIST}\index{LIST}\index{BASIC 65-commando's!LIST}-commando.


\subsection{Oefeningen om te proberen}

{\bf 1. Kunt u het tot een hoger of lager getal laten tellen?}

Op dit moment telt het van 1 tot 10. Kunt u het veranderen zodat het tot 20 telt? Of zodat het van 3 tot 17 telt? Of van 14.5 tot 21.5? Wat denkt u nodig te hebben om de telvolgorde om te keren?

{\em Aanwijzing:} U zult het \screentextwide{FOR}-commando moeten aanpassen.

{\bf 2. Kunt u de telstap veranderen?}

Op dit moment telt het met stapjes van één, d.w.z. elk getal is één meer dan het vorige. Kunt u het veranderen zodat het met stapjes van twee telt? Of met halve stapjes, zodat het 1, 1.5, 2, 2.5, 3, \ldots telt?

{\em Aanwijzing:} U zult de \screentextwide{STEP}-clausule van het \screentextwide{FOR}-commando moeten aanpassen.\index{STEP}\index{BASIC 65-commando's!STEP}


{\bf 3. Kunt u één van de tafels van vermenigvuldiging laten weergeven?}

Op dit moment geeft het de antwoorden van de tafel van 1 weer, omdat het met stapjes van één telt. Kunt u het met stapjes van drie laten tellen en de tafel van 3 weergeven?

{\em Aanwijzing:} U zult het \screentextwide{FOR}-commando moeten aanpassen.

{\bf 4. Kunt u de tafels van vermenigvuldiging van 1$\times$1 tot 10$\times$10 laten weergeven?}

{\em Aanwijzing:} U kunt een puntkomma (;) achter het \screentextwide{PRINT}-commando zetten, zodat er meer dan één waarde per regel op het scherm past.\\  {\em Aanwijzing:} Het \screentextwide{PRINT}-commando zonder argument springt gewoon naar het begin van de volgende regel.\\  {\em Aanwijzing:} Mogelijk heeft u meerdere \screentextwide{FOR}-lussen nodig, de één binnen de andere.

\section{Eerste stappen met tekst en getallen}

In het vorige gedeelte gebruikten we zowel getallen als tekst. Tekst in een computer ontstaat door afzonderlijke letters en andere symbolen aaneen te rijgen. Zo'n aaneenschakeling noemen we {\em strings}. Ook de afzonderlijke letters en symbolen hebben een naam: we noemen ze {\em tekens}. Het Engelse woord character komt uit de drukkerswereld, waar hij verwijst naar elk symbool dat op een bladzijde gedrukt kan worden. Voor computers heeft het vrijwel dezelfde betekenis en de verzameling tekens die een computer kan weergeven noemen we, nogal fantasieloos, een {\em tekenset}.\index{teken}\index{tekenset}\index{string}.

Indien de MEGA65 een vorm van invoer verwacht, zoekt hij gewoonlijk naar één van vier dingen:

\begin{enumerate}
\item {\em een sleutelwoord} zoals \screentextwide{PRINT} of \screentextwide{STEP}, woorden die een speciale betekenis voor de computer hebben;\item {\em een variabelenaam} zoals \screentextwide{I} of \screentextwide{A\$}, welke de computer vervolgens gebruikt om een stukje informatie op te slaan of op te halen;\item {\em een getal} zoals \screentextwide{42} of \screentextwide{-30.3137}; of\item {\em een string} zoals \screentextwide{"HALLO COMPUTER"} of \screentextwide{"23 KILOMETER"}.
\end{enumerate}

\needspace{4cm} % Dont allow following paragraph to separate from
                % following screenshot
Soms kunt u kiezen welk soort invoer u geeft, terwijl die keuze andere keren beperkter is. Wat de computer aanvaardt, hangt af van wat u op dat moment doet. In het vorige gedeelte ontdekten we bijvoorbeeld dat u de computer een sleutelwoord of een getal kunt geven wanneer hij \screentextwide{READY} weergeeft. Denkt u dat de computer alle vier de soorten aanvaardt wanneer hij \screentextwide{READY.} weergeeft? We weten al dat u sleutelwoorden en getallen kunt invoeren, maar hoe zit het met variabelenamen en strings? Laten we een variabelenaam intypen, bijvoorbeeld \screentextwide{N}, \specialkey{RETURN} indrukken en kijken wat er gebeurt. Daarna proberen we het met een string, bijvoorbeeld \screentextwide{"DIT IS EEN STRING"}.

\screenshotwrap{dutch/images/getting-started/typing-variable-name-or-string}

Elke keer krijgt u een syntaxisfout, welke u vertelt dat de computer de invoer die u hem gegeven heeft niet begrijpt. Laten we beginnen met de omstandigheid waarin u de variabele typte: als u de computer alleen maar de naam van een opgeslagen stukje gegevens geeft, heeft hij er geen flauw benul van wat u wilt dat hij doet. Hetzelfde geldt wanneer u de computer een stukje gegevens geeft, zoals een string, zonder hem te vertellen wat hij ermee moet doen.

Zoals we in het vorige gedeelte echter hebben gezien, kunnen we de computer met het \screentextwide{PRINT}-commando vertellen dat we het stukje gegevens dat in een variabele is opgeslagen, willen zien. In plaats daarvan zouden we dus \screentextwide{PRINT N} kunnen intypen; de computer weet dan wat hij moet doen en geeft het stukje gegevens weer dat in de variabele \screentextwide{N} is opgeslagen.

Het gebruik van het \screentextwide{PRINT}-commando is zelfs zo algemeen dat programmeurs het vervelend vonden om steeds het \screentextwide{PRINT}-commando te moeten intypen. Ze bedachten daarom een sneltoets: typt u een vraagteken, d.w.z. een \screentextwide{?}, dan weet de computer dat u \screentextwide{PRINT} bedoelt. Typt u bijvoorbeeld \screentextwide{? N}, dan gebeurt hetzelfde als wanneer u \screentextwide{PRINT N} typt. U moet na elk commando uiteraard \specialkey{RETURN} indrukken om de computer te vertellen dat u wilt dat hij verwerkt wat u getypt heeft. Vanaf nu gaan we ervan uit dat u daaraan denkt zonder dat we u eraan hoeven te herinneren.

\needspace{4cm} % Dont allow following paragraph to separate from
                % following screenshot
De \screentextwide{?}-sneltoets werkt ook wanneer u de computer een commando als onderdeel van een programma laat onthouden. Typt u dus \screentextwide{1 ? N} en daarna \screentextwide{LIST}, dan ziet u \screentextwide{1 PRINT N}, zoals te zien is in de volgende schermafdruk:

\screenshotwrap{dutch/images/getting-started/print-question-mark}

Zoals we in het vorige gedeelte zagen, is er nog geen waarde in de variabele \screentextwide{N} opgeslagen; wanneer de computer dus kijkt wat erin staat, vindt hij niets. Omdat \screentextwide{N} een {\em numerieke variabele}\index{variabele!numeriek} is, betekent ``niets'' hier nul. Was het een {\em stringvariabele}\index{variabele!string}, dan had hij letterlijk niets gevonden. Dat kunnen we uitproberen, maar eerst moeten we uitleggen hoe we de computer vertellen dat we het over een stringvariabele hebben. Dat doen we door een dollarteken, d.w.z. een \screentextwide{\$}, achter de naam van de variabele te zetten. Zetten we dus een \screentextwide{\$} achter de variabelenaam \screentextwide{N}, dan verwijst dat naar een stringvariabele genaamd \screentextwide{N\$}.

\needspace{4cm} % Dont allow following paragraph to separate from
                % following screenshot
Met het inmiddels vertrouwde \screentextwide{PRINT}-commando (of de \screentextwide{?}-sneltoets) kunnen we met deze variabelen experimenteren en bekijken wat erin staat. We hebben echter een eenvoudige manier nodig om er waarden in te zetten. Gelukkig zijn we niet de eersten die waarden in variabelen willen zetten; daarom bestaat er een \screentextwide{LET}\index{LET}\index{BASIC 65-commando's!LET}-commando. Met het \screentextwide{LET}-commando zet u een waarde in een variabele. Zo kunnen we de computer bijvoorbeeld vertellen:

\begin{screencode}
  LET N = 5.3
\end{screencode}

\needspace{4cm} % Dont allow following paragraph to separate from
                % following screenshot
Hiermee vertelt u de computer dat hij de waarde 5.3 in de variabele \screentextwide{N} moet zetten. Met het \screentextwide{PRINT}-commando kunnen we vervolgens controleren of dat gelukt is. Op dezelfde manier kunnen we een waarde in de variabele \screentextwide{N\$} zetten, bijvoorbeeld met een regel als:

\begin{screencode}
  LET N$ = "DE KONING DER AARDAPPELMENSEN"
\end{screencode}

\needspace{4cm} % Dont allow following paragraph to separate from
                % following screenshot
Als we deze uitproberen, zien we ongeveer het volgende:

\screenshotwrap{dutch/images/getting-started/let-command-examples}

\needspace{4cm} % Dont allow following paragraph to separate from
                % following screenshot
Eerder vermeldden we al dat \screentextwide{N} een numerieke variabele is en \screentextwide{N\$} een stringvariabele. Dat betekent dat we alleen getallen in \screentextwide{N} kunnen zetten en alleen strings in \screentextwide{N\$}. Proberen we het verkeerde soort gegevens in een variabele te zetten, dan laat de computer ons weten dat het soort gegevens niet past op de plaats waar we het willen zetten. Hij geeft ons dan een \screentextwide{TYPE MISMATCH  ERROR}\index{Fouten!Type mismatch}\index{Type mismatch error}, zoals hier:

\screenshotwrap{dutch/images/getting-started/type-mismatch-errors}

Dit brengt ons bij een vrij belangrijk punt: \screentextwide{N} en \screentextwide{N\$} zijn afzonderlijke variabelen, ook al lijken hun namen op elkaar. Dit geldt voor alle mogelijke variabelenamen: eindigt de variabelenaam op een \screentextwide{\$}-teken, dan is het een stringvariabele die geheel losstaat van de gelijknamige numerieke variabele. Om een beetje jargon te gebruiken: elke variabele van een bepaald {\em type} heeft zijn eigen afzonderlijke {\em naamruimten}\index{naamruimten}.

(Er zijn overigens nog vier andere variabelenaamruimten die we nog niet hebben besproken: integer-variabelen, herkenbaar aan een \screentextwide{\%}-teken achter de naam, bijv. \screentextwide{N\%} en arrays van numerieke variabelen, stringvariabelen of integer-variabelen. Maakt u zich daar voorlopig nog geen zorgen over; daar komen we later nog op terug.)

Tot nu toe hebben we variabelen alleen waarden gegeven in de directe modus, of met behulp van constructies als \screentextwide{FOR}-lussen. We hebben echter nog niet gezien hoe we gegevens van de gebruiker kunnen krijgen terwijl een programma draait. Een manier om dat te doen is met het
\screentextwide{INPUT}\index{BASIC 65-commando's!INPUT}\index{INPUT}-commando.

\needspace{4cm} % Dont allow following paragraph to separate from
                % following screenshot
\screentextwide{INPUT} is heel eenvoudig te gebruiken: we hoeven alleen maar aan te geven in welke variabele de invoer moet komen. Om de computer bijvoorbeeld de gebruiker te laten vragen iets in te typen dat in de variabele \screentextwide{A\$} moet komen, kunnen we zoiets gebruiken als \screentextwide{INPUT A\$}. Het enige nadeel van het \stw{INPUT}-commando is dat u het niet in de directe modus\index{Directe modus} kunt gebruiken. Probeert u het toch, dan meldt de computer \stw{ILLEGAL DIRECT  ERROR}\index{Fouten!Illegal Direct}\index{Illegal Direct Error}. Probeer het maar; u zult ongeveer het volgende zien:

\screenshotwrap{images/getting-started/illegal-direct-error}

\needspace{4cm}
Dat betekent dat het \stw{INPUT}-commando alleen als onderdeel van een programma gebruikt kan worden. In plaats daarvan kunnen we dus iets doen als volgt:

\begin{screencode}
1 INPUT A$
2 PRINT "U VOERDE IN "; A$
RUN
\end{screencode}

\needspace{4cm}
Wat denkt u dat dit zal doen? De eerste regel vraagt de computer om invoer die in de variabele \stw{A\$} gezet moet worden en de tweede regel geeft de string \stw{"U VOERDE IN"} weer, gevolgd door wat het \stw{INPUT}-commando van de gebruiker heeft gelezen. Laten we het uitproberen:

\screenshotwrap{dutch/images/getting-started/input-example-1}

Had u dit verwacht? Wat doet dat vraagteken daar? Het \stw{?} hier is de wijze waarop de computer u laat weten dat een programma op invoer van u wacht. Dat betekent dat de computer hetzelfde teken, \stw{?}, gebruikt om twee verschillende dingen aan te duiden: typt u het als onderdeel van een programma of in de directe modus, dan is het een verkorte vorm van het \stw{PRINT}-commando. Dat is zo wanneer u het intypt. Toont de computer het evenwel aan u, dan heeft het een andere betekenis: namelijk dat de computer wacht tot u iets intypt. Er is ook nog een derde wijze waarop het \stw{?}-teken gebruikt wordt. Heeft u gemerkt hoe? Het geeft het begin van een foutmelding aan. Een syntaxisfout wordt bijvoorbeeld aangegeven met \stw{?SYNTAX ERROR}. Indien een teken of iets anders in verschillende omstandigheden of contexten verschillende betekenissen heeft, zeggen we dat het {\em contextafhankelijk}\index{contextafhankelijk} is.

\needspace{4cm}
Laten we terugkeren naar ons voorbeeld: als u nu iets intypt en op \specialkey{RETURN} drukt om de computer te laten weten dat u klaar bent, gaat het programma verder, zoals hier:

\screenshotwrap{dutch/images/getting-started/input-example-2}

\needspace{4cm}
Natuurlijk wisten we niet echt wat we moesten intypen, want het programma gaf de gebruiker geen enkele aanwijzing over wat de programmeur wilde dat hij deed. We zouden dus moeten proberen wat instructies te geven. Als we de gebruiker bijvoorbeeld zijn naam wilden laten typen, konden we een boodschap weergeven met het verzoek zijn naam te typen, zoals hier:

\begin{screencode}
  1 PRINT "HOE HEET U"
  2 INPUT A$
  3 PRINT "HALLO "; A$
\end{screencode}

\needspace{4cm}
Als we dit programma nu uitvoeren, krijgt de gebruiker een aanwijzing over wat we van hem verwachten en de hele ervaring zal voor hem een stuk duidelijker worden:

\screenshotwrap{dutch/images/getting-started/input-example-3}

Wanneer we het programma uitvoeren, zien we eerst de boodschap \stw{HOE HEET U} uit regel 1. De computer geeft de dubbele aanhalingstekens niet weer, omdat die de computer alleen maar vertelden dat het stukje gegevens ertussen een string is. De string zelf is alleen het deel dat ertussen staat.

Daarna zien we opnieuw het \stw{?}-teken en de knipperende cursor, welke ons vertelt dat de computer op invoer van ons wacht. De rest van het programma is {\em geblokkeerd}\index{I/O!blokkeren}\index{geblokkeerd} en kan niet verdergaan totdat we het stukje gegevens hebben getypt. Zodra we de invoer hebben getypt, slaat de computer die op in de variabele \stw{A\$} en kan hij verder. Wanneer hij dus bij regel 3 van het programma komt, heeft hij alles wat hij nodig heeft en geeft hij zowel de \stw{HALLO}-boodschap weer als de gegevens die in de variabele \stw{A\$} zijn opgeslagen.

Merk op dat het woord \stw{LENNART} nergens in het programma voorkomt. Het bestaat alleen in de variabele. Dit vermogen om gegevens te verwerken die geen deel van het programma uitmaken, is een van de dingen die computerprogramma's zo krachtig en geschikt maken voor zoveel doeleinden. We hoeven alleen maar de invoer te wijzigen en we kunnen dan andere uitvoer krijgen.


\needspace{4cm}
Ons programma kunnen we bijvoorbeeld keer op keer uitvoeren en elke keer andere invoer geven, waarna het programma zijn uitvoer aanpast aan wat we typen. Vrij handig, nietwaar? Laat de rest van de bemanning het maar eens uitproberen:

\screenshotwrap{dutch/images/getting-started/input-extra-ignored-1}

We kunnen zien dat het programma elke keer de boodschap weergeeft, aangepast aan de invoer die u getypt heeft\ldots\  Tot we bij \stw{REMWAALS, HEER} komen. Zoals altijd veroorzaakt meneer Remwaals problemen. In dit geval kon hij het niet laten om zijn aanspreekvorm achter zijn naam te zetten.

De computer heeft ons een soort foutmelding gegeven, \stw{?EXTRA IGNORED}\index{Extra  Ignored}\index{Fouten!Extra Ignored}\index{Waarschuwingen!Extra Ignored}. De melding is niet in het rood geschreven en het woord \stw{ERROR} ontbreekt op het eind. Dat betekent dat het een waarschuwing is in plaats van een fout. Omdat het slechts een waarschuwing is, gaat het programma gewoon verder. Er is echter wel iets gebeurd: de computer heeft de aanspreekvorm \stw{HEER} van meneer Remwaals genegeerd, dat wil zeggen dat hij de extra invoer heeft genegeerd. Dat komt doordat het \stw{INPUT}-commando niet daadwerkelijk een hele regel invoer leest. Het leest veeleer {\em één stukje gegevens}. Het \stw{INPUT}-commando neemt aan dat een stukje gegevens eindigt aan het eind van een invoerregel of wanneer het een komma (\stw{,}) of dubbele punt (\stw{:}) tegenkomt.\index{, (komma)}\index{: (dubbele punt)}

\needspace{4cm}
Wilt u een van die symbolen gebruiken, dan moet u het hele stukje gegevens tussen dubbele aanhalingstekens zetten. Als meneer Remwaals deze handleiding had gelezen in plaats van de computer op de bonnefooi te gebruiken, had hij geweten dat hij \stw{"REMWAALS, HEER"} moest typen, compleet met de dubbele aanhalingstekens, om het programma goed te laten draaien. Het is belangrijk dat de aanhalingstekens om het hele stukje gegevens staan, want anders denkt de computer dat het eerste aanhalingsteken het begin van een nieuw stukje gegevens aangeeft. Hieronder ziet u het verschil dat dit maakt:

  \screenshotwrap{dutch/images/getting-started/input-quoting-1}

\needspace{1.5cm}
Al kan dit soms behoorlijk vervelend zijn, het heeft een doel: met het \stw{INPUT}-commando kunt u meer dan één stukje gegevens lezen. Dat doet u door meer dan één variabele achter het \stw{INPUT}-commando te zetten, telkens gescheiden door een komma. Het \stw{INPUT}-commando verwacht dan meerdere stukjes gegevens. Zo zouden we bijvoorbeeld kunnen vragen om iemands naam en leeftijd, met een programma als dit:

\begin{screencode}
  1 PRINT "HOE HEET U EN WAT IS UW LEEFTIJD"
  2 INPUT A$, A
  3 PRINT "HALLO "; A$
  4 PRINT "U BENT"; A; " JAAR OUD."
\end{screencode}

\needspace{4cm}
Voeren we dit programma uit, dan kunnen we de twee stukjes gegevens op één regel typen wanneer de computer ons de \stw{?}-prompt toont, bijvoorbeeld \stw{LENNART, 3000000}. Let op de komma die de twee stukjes gegevens, \stw{LENNART} en \stw{3000000}, van elkaar scheidt. Het is ook de moeite waard om op te merken dat we in het getal 3000000 geen scheidingstekens voor duizendtallen hebben gezet. Hadden we dat wel gedaan, dan zou de computer denken dat we drie afzonderlijke stukjes gegevens bedoelden, \stw{3}, \stw{000} en \stw{000} en dat is niet wat we bedoelden. Laten we dus eens bekijken hoe het eruitziet wanneer we \stw{LENNART, 3000000} als invoer aan het programma geven:

  \screenshotwrap{dutch/images/getting-started/input-multiple-1}

In dit geval leest het \stw{INPUT}\index{INPUT}\index{BASIC 65-commando's!INPUT}-commando de twee stukjes gegevens en zet het eerste in de variabele \stw{A\$} en het tweede in de variabele \stw{A}. Wanneer het programma bij regel 3 komt, geeft het \stw{HALLO} weer, gevolgd door het eerste stukje gegevens. Wanneer het vervolgens bij regel 4 komt, geeft het de string \stw{U BENT} weer, gevolgd door de inhoud van de variabele \stw{A}, het getal 3000000 en ten slotte de string \stw{JAAR OUD}.

Het is ook mogelijk om telkens maar één stukje gegevens tegelijk te geven. In dat geval vraagt het \stw{INPUT}-commando met een prompt van twee vraagtekens om het tweede gegeven, d.w.z. \stw{??}. Zodra het het tweede gegeven heeft.  (Staan er meer dan twee variabelen op het \stw{INPUT}-commando, dan toont de computer nog steeds dezelfde \stw{??}-prompt, in plaats van steeds meer vraagtekens af te drukken.)

\needspace{4cm}
Als we dit met ons programma proberen, zien we deze \stw{?}- en \stw{??}-prompts en hoe het eerste stukje gegevens in \stw{A\$} terechtkomt, omdat dit de eerste variabele in het \stw{INPUT}-commando is. Het tweede gegeven komt in \stw{A} terecht, omdat \stw{A} de tweede variabele na het \stw{INPUT}-commando is. Zo ziet het eruit wanneer we deze invoer aan ons programma geven:


\screenshotwrap{dutch/images/getting-started/input-multiple-2}

Tot nu toe hebben we de gebruiker om gegevens gevraagd door met een \stw{PRINT}-commando de boodschap weer te geven en vervolgens met een \stw{INPUT}-commando aan de computer kenbaar te maken in welke variabelen we de gegevens willen opnemen. Net als bij het \stw{PRINT}-commando gebeurt dit echter zo vaak dat er een verkorte vorm voor bestaat. Het heeft bovendien het voordeel dat het er tijdens het uitvoeren netter uitziet en het programma iets korter maakt. De verkorte vorm bestaat eruit de te tonen boodschap na het \stw{INPUT}-commando te plaatsen, maar vóór de eerste variabele.

We kunnen ons programma aanpassen om deze werkwijze te gebruiken. Eerst kunnen we regel 3 zodanig wijzigen dat de prompt na het \stw{INPUT}-commando komt. Dat kunnen we op twee manieren doen. De eerste is eenvoudig een nieuwe regel 3 in te typen. De computer vervangt dan automatisch de oude regel 3 door de nieuwe.

Zoals we inmiddels al een paar keer hebben vermeld, zijn programmeurs luie wezens en dus is er een verkorte vorm: als u de regel die u wilt wijzigen op het scherm kunt zien, kunt u met de cursortoetsen naar die regel gaan, hem op het scherm bewerken en vervolgens op \specialkey{RETURN} drukken om de computer de nieuwe versie van de regel te laten aanvaarden.\index{Programma's!bewerken}\index{Programma's!regels vervangen}\index{Regels!bewerken}\index{Regels!vervangen}

\needspace{4cm}
Hoe dan ook, u kunt controleren of de wijzigingen zijn geslaagd door het \stw{LIST}-commando op een willekeurige lege regel van het scherm in te typen. Dit toont de herziene versie van het programma. Bijvoorbeeld:

\screenshotwrap{dutch/images/getting-started/replacing-line-1}

\needspace{3cm}
Toch hebben we nog een klein probleempje: regel 1 geeft de boodschap \stw{HOE HEET U EN WAT IS UW LEEFTIJD} weer, waarna regel 2 hem nog eens weergeeft! We willen dat de boodschap slechts één keer verschijnt. We willen daarom regel 1 zodanig wijzigen dat hij dit niet meer doet. Omdat er op regel 1 geen ander commando staat dat we willen behouden, kan die regel gewoon leeg worden. We kunnen dus zoiets intypen:

\begin{screencode}
1
\end{screencode}

\needspace{4cm}
We kunnen bevestigen dat de inhoud van de regel is verwijderd door het \stw{LIST}-commando opnieuw uit te voeren, zoals hier:

\screenshotwrap{dutch/images/getting-started/deleting-line-1}

Heeft u iets interessants opgemerkt? Toen we de computer opdroegen regel 1 van het programma leeg te maken, heeft hij hem volledig verwijderd! Dat komt doordat de computer van mening is dat een lege regel nergens voor dient. Hij zorgt er bovendien voor dat uw programma's niet vol lege regels komen te staan wanneer u veel wijzigingen in uw programma's aanbrengt.

Het is ook mogelijk om met DELETE een reeks regels te verwijderen. U kunt bijvoorbeeld (maar doe dit nu nog niet) regels 3-4 verwijderen met:

\begin{screencode}
  DELETE 3-4
\end{screencode}

U kunt meer lezen over het DELETE-commando in de Naslag BASIC 65-commando's.

\needspace{4cm}
Nu we dat uit de weg hebben, kunt u ons programma nu uitvoeren en kijken wat er gebeurt. Typ zoals gebruikelijk gewoon het \stw{RUN}-commando in en druk op \specialkey{RETURN}. U ziet dan zoiets:

\screenshotwrap{dutch/images/getting-started/input-comma}

\needspace{2.5cm}
We zien daar onze prompt \stw{HOE HEET U EN WAT IS UW LEEFTIJD}, maar nu verschijnt de cursor zonder enig vraagteken (\stw{?}). Dat komt doordat we een komma (\stw{,}) achter de boodschap in het \stw{INPUT}-commando hebben gezet. Om het vraagteken te krijgen moeten we in plaats daarvan een puntkomma (\stw{;}) achter de boodschap zetten, zoals dit:

\begin{screencode}
INPUT "HOE HEET U EN WAT IS UW LEEFTIJD"; A$, A
\end{screencode}

\needspace{4cm}
Als u het programma nu uitvoert, zou u moeten zien wat we zoeken:

\screenshotwrap{dutch/images/getting-started/input-semicolon}

  \subsection{Oefeningen om te proberen}

{\bf 1. Kunt u het programma iemand naar zijn of haar naam laten vragen en vervolgens naar de lievelingskleur?}

Op dit moment vraagt het programma naar de naam en de leeftijd. Kunt u het zo veranderen dat het in plaats van de leeftijd de lievelingskleur weergeeft?

{\em Aanwijzing:} Wat voor soort gegevens is de leeftijd? Is het een numerieke waarde of een string? Is dat hetzelfde soort gegevens als de naam van een kleur?

{\bf 2. Kunt u een programma schrijven dat iemand naar zijn of haar naam vraagt, de begroeting weergeeft en vervolgens naar de leeftijd vraagt en dat antwoord weergeeft?}

  \needspace{2cm}
Op dit moment verwacht het programma beide stukjes gegevens tegelijkertijd. Dit betekent dat het programma geen bericht over het eerste stukje gegevens kan weergeven voordat het beide stukjes heeft ontvangen. Verander het programma zodat u in plaats daarvan een interactie krijgt zoals de volgende:

\begin{screencode}
HOE HEET U? DIEPE GEDACHTE
HALLO DIEPE GEDACHTE
WAT IS HET ANTWOORD? 42
VOLGENS U WAS HET ANTWOORD 42
\end{screencode}

{\em Aanwijzing:} U zult meer regels nodig hebben in uw programma, zodat u meer dan één \stw{INPUT}- en \stw{PRINT}-commando kunt gebruiken.

{\bf 3. Kunt u een programma schrijven dat enkele vragen stelt en vervolgens de lijst met gegeven antwoorden weergeeft?}

\needspace{2cm}
Bedenk enkele vragen die u iemand zou willen stellen en schrijf vervolgens een programma dat die vragen stelt, de antwoorden onthoudt en ze met een passende melding weergeeft. Het uitvoeren van uw programma zou er bijvoorbeeld als volgt uit kunnen zien:

\begin{screencode}
  HOE HEET U? FRODO
  HOE OUD BENT U? 33
  WAT IS UW LIEVELINGSETEN? ALLES!
  BEDANKT VOOR DE ANTWOORDEN.
  U HEET FRODO
  U BENT 33 JAAR OUD
  UW LIEVELINGSETEN IS ALLES!
\end{screencode}

{\em Aanwijzing:} U zult meer regels in uw programma nodig hebben om de verschillende \stw{INPUT}- en \stw{PRINT}-commando's op te nemen.

{\em Aanwijzing:} U zult goed moeten nadenken over welke variabelenamen u zult gebruiken.

\section{Eenvoudige beslissingen nemen}

In het vorige gedeelte hebben we geleerd hoe we tekst en numerieke gegevens invoeren en weergeven. De programma's hebben echter steeds gewoon de instructieregels in volgorde gevolgd, zonder dat ze op basis van de ingevoerde gegevens konden beslissen wat er moest gebeuren.

In dit gedeelte zien we hoe we eenvoudige beslissingen kunnen nemen met behulp van de \stw{IF}\index{IF}\index{BASIC 65-commando's!IF}- en \stw{THEN}\index{THEN}\index{BASIC 65-commando's!THEN}-commando's. Het \stw{IF}-commando controleert of iets waar of onwaar is en als dat het geval is, zorgt het ervoor dat de computer het commando uitvoert dat na het \stw{THEN}-commando staat.

De computer beslist of iets waar of onwaar is door de aangeleverde gegevens te bewerken met behulp van één van meerdere symbolen. Deze symbolen worden daarom {\em operatoren} genoemd. Omdat ze twee dingen met elkaar vergelijken, hangen ze bovendien af van de relatie tussen die dingen. Om die reden worden ze {\em relationele} operatoren genoemd.\index{vergelijkingsoperatoren}\index{vergelijkingsoperatoren}Daartoe behoren de volgende:

\begin{itemize}
\item Gelijkheid (\stw{=}). Zo zou \stw{3 = 3} waar zijn, terwijl \stw{3 = 2} onwaar zou zijn.  \item Kleiner dan (\stw{<}). Zo zou \stw{1 < 3} waar zijn, terwijl zowel \stw{3 < 3} als \stw{1 < 3} onwaar zou zijn.  \item Groter dan (\stw{>}). Zo zou \stw{3 > 1} waar zijn, terwijl zowel \stw{3 > 3} als \stw{1 > 3} onwaar zou zijn.
\end{itemize}

Omdat het vaak voorkomt dat u wilt nagaan of iets gelijk is aan of groter dan iets anders, of gelijk is aan of kleiner dan iets anders, bestaan hiervoor verkorte schrijfwijzen. Wilt u daarnaast testen of iets niet gelijk is aan iets anders, dan is daar ook een relationele operator voor:

\begin{itemize}
\item Ongelijkheid, in het Engels gewoonlijk {\em not equal} genoemd\index{niet gelijk aan}\index{ongelijkheid}\index{<> (ongelijk aan)} (\stw{<>}). Dit wijkt af van het wiskundige teken voor niet-gelijk, $\ne$, omdat de tekenset van de MEGA65 geen teken bevat dat daarop lijkt. De programmeurs die BASIC voor de MEGA65 schreven, gebruikten daarom de tekens voor groter dan en kleiner dan samen om aan te geven dat iets kleiner dan of groter dan is, dat wil zeggen: niet gelijk aan. Zo zou \stw{1 <> 3} waar zijn, terwijl \stw{3 <> 3} onwaar zou zijn.  \item Kleiner dan of gelijk aan (\stw{<=}). Zo zouden \stw{1 < 3} en \stw{3 <= 3} waar zijn, terwijl \stw{4 < 3} onwaar zou zijn.  \item Groter dan of gelijk aan (\stw{>=}). Zo zouden \stw{3 >= 1} en \stw{3 >= 3} waar zijn, terwijl \stw{1 >= 3} onwaar zou zijn.
\end{itemize}

Een handig trucje wanneer u zich niet meer kunt herinneren hoe de tekens \stw(<) en \stw(>) wijzen: de kant met de meeste lijnuiteinden is de kant waar het grootste getal moet staan. Zo heeft het teken \stw(<) links één punt, maar aan de rechterkant twee lijnuiteinden. Wil een uitspraak met \stw(<) waar zijn, dan moet het getal aan de linkerkant kleiner zijn dan het getal aan de rechterkant. Dit trucje werkt zelfs voor het gelijkteken, \stw(=), want dat heeft aan beide kanten evenveel uiteinden; u kunt dus onthouden dat de getallen aan beide kanten gelijk moeten zijn. Het werkt ook wanneer u twee tekens samen gebruikt, zoals \stw(>=): het is waar wanneer de voorwaarde voor één van de tekens waar is. In dit geval heeft het teken \stw(>) dus meer uiteinden aan de linkerkant dan aan de rechterkant, zodat het waar is wanneer het getal aan de linkerkant groter is dan het getal aan de rechterkant. Aangezien het teken \stw(=) aan elke kant twee uiteinden heeft, is het ook waar wanneer de twee getallen gelijk zijn.

\needspace{2cm}
Met deze relationele operatoren kunnen we een regel schrijven die iets doet, maar alleen wanneer iets waar of onwaar is. U kunt dit eens uitproberen met een paar voorbeelden:

\begin{screencode}
  IF -2 < 0 THEN PRINT "-2 IS EEN NEGATIEF GETAL"
  IF 2 < 0 THEN PRINT "2 IS A EEN NEGATIEF GETAL"
  IF 0 < -2 THEN PRINT "-2 IS EEN POSITIEF GETAL"
  IF 0 < 2 THEN PRINT "2 IS EEN POSITIEF GETAL"
\end{screencode}

\needspace{4cm}
Deze commando's werken prima in de directe modus, dus u kunt ze zo direct in de computer intypen om te zien wat ze doen. Dat kan handig zijn om te testen of u de logica goed heeft wanneer u een \stw{IF} -- \stw{THEN}-commando opstelt. Typt u deze commando's in, dan ziet u ongeveer het volgende:

\screenshotwrap{dutch/images/getting-started/if-then-less-than-examples}

We kunnen zien dat alleen de \stw{PRINT}-commando's werden uitgevoerd die op een \stw{IF}-commando volgden waarvan de waarde waar was. De computer negeerde de rest zwijgend. We kunnen deze commando's echter ook in een programma opnemen. Laten we dus een programmaatje maken dat om twee getallen vraagt en zegt of ze gelijk zijn, of dat het ene groter of kleiner is dan het andere. Denk, voordat u naar het programma kijkt, eerst na over hoe u het zou aanpakken en kijk of u er zelf uit kunt komen. De aanwijzing die wij u geven, is dat het \stw{IF}-commando ook de naam van een variabele aanvaardt, niet alleen getallen. U kunt dus bijvoorbeeld iets doen als \stw{IF A  > B THEN PRINT "IETS"}. Het programma staat op de volgende bladzijde, zodat u niet kunt spieken voordat u erover nagedacht heeft!

\pagebreak

Heeft u het geprobeerd? Er zijn veel verschillende manieren waarop het gedaan kan worden, maar dit is wat wij bedacht hebben:

\begin{screencode}
1 INPUT "WAT IS HET EERSTE GETAL"; A
2 INPUT "WAT IS HET TWEEDE GETAL"; B
3 IF A = B THEN PRINT "DE GETALLEN ZIJN GELIJK"
4 IF A > B THEN PRINT "HET EERSTE GETAL IS GROTER"
5 IF B > A THEN PRINT "HET TWEEDE GETAL IS GROTER"
\end{screencode}

Vervolgens kunnen we het programma zo vaak uitvoeren als we willen en de computer kan ons vertellen welke van de twee getallen die we hem geven het grootste is, of dat ze gelijk zijn:

\screenshotwrap{dutch/images/getting-started/if-compare-variables-1}

Let u erop dat we in dit programma geen vaste getallen in het \stw{IF}-commando gebruikten, maar in plaats daarvan namen van variabelen. Dit is een van de krachtigste eigenschappen van het programmeren, samen met de mogelijkheid om op basis van gegevens beslissingen te nemen. Doordat de programmeur gegevens bij naam kan aanspreken, ongeacht hun huidige waarde of hoe ze daar gekomen zijn, kan hij zeer flexibele programma's maken.

Laten we eens nadenken over een iets interessanter voorbeeld: een spelletje ``raad het getal''.\index{Raad het getal}\index{Spellen!Raad het getal} Daarvoor hebben we een getal nodig dat iemand moet raden en vervolgens moeten we de gokken aannemen en aangeven of de gok juist was of niet. Is de gok onjuist, dan moeten we de gebruiker vertellen of het juiste getal hoger of lager is.

De meeste ingrediënten voor zo'n programma hebben we al geleerd: we kunnen \stw{LET} gebruiken om een variabele op het geheime getal te zetten, \stw{INPUT} om de gebruiker naar zijn gok te vragen en vervolgens \stw{IF}, \stw{THEN} en \stw{PRINT} om de gebruiker te vertellen of zijn gok juist was of niet. Laten we dus iets maken. Ook nu geldt: als u wilt, kunt u een paar minuten stilstaan, nadenken en experimenteren om te zien of u zo'n programma zelf kunt maken.

Zo hebben wij het gedaan. U hoeft zich echter geen zorgen te maken als u het op een heel andere manier heeft gedaan: er zijn vaak veel manieren om een programma te schrijven dat een bepaalde taak uitvoert.

\begin{screencode}
1 SN=23
2 PRINT "RAAD EEN GETAL TUSSEN 1 EN 100"
3 INPUT "WAT IS UW GOK"; G
4 IF G<SN THEN PRINT "MIJN GETAL IS GROTER"
5 IF G>SN THEN PRINT "MIJN GETAL IS KLEINER"
6 IF G=SN THEN PRINT "GEFELICITEERD! U HEEFT MIJN GETAL GERADEN!"
\end{screencode}

\needspace{4cm}
De eerste regel zet ons geheime getal in de variabele \stw{SN}. De tweede regel toont een bericht dat de gebruiker vertelt wat de bedoeling is. De derde regel vraagt de gebruiker naar zijn gok en zet die in de variabele \stw{G}. De vierde, vijfde en zesde regel controleren vervolgens of de gok juist is of niet en zo niet, welk bericht er getoond moet worden. Dat gebeurt met het \stw{IF}-commando en een passende relationele operator bij elke beslissing. Dit werkt goed, maar slechts tot op zekere hoogte. Bijvoorbeeld:

\screenshotwrap{dutch/images/getting-started/guess-number-1}

We kunnen zien dat het programma het bericht toont, om een gok vraagt en gepast reageert. Willen we echter nog eens raden, dan moeten we voor elke extra gok opnieuw het \stw{RUN}-commando gebruiken. Dat is vanuit het perspectief van de gebruiker niet zo mooi. Toch zal dat waarschijnlijk niet lang een probleem zijn, want de gebruiker kan het geheime getal in de programmalijst op het scherm zien!

Deze problemen willen we dus graag oplossen. Laten we beginnen met het verbergen van de programmalijst. Eerder noemden we al dat wanneer het scherm scrollt, alles wat bovenaan het scherm stond verdwijnt. We kunnen er dus simpelweg voor zorgen dat het scherm voldoende scrollt, zodat een zichtbare programmalijst niet langer zichtbaar is. Dat kunnen we doen met \stw{PRINT} en een \stw{FOR}-lus. Het scherm telt 25 regels, dus we kunnen iets doen als:

\begin{screencode}
FOR I = 1 to 25
PRINT
NEXT I
\end{screencode}

Dit kan evenwel op een betere manier aangepakt worden. Als u \specialkey{SHIFT} ingedrukt houdt en vervolgens \specialkey{CLR\\HOME} indrukt, wordt het scherm gewist. Dat is veel eenvoudiger en handiger. Hoe kunnen we zoiets echter in ons programma voor elkaar krijgen? Dat blijkt heel eenvoudig: u kunt het typen terwijl u een string invoert! Het toetsenbord werkt namelijk anders, afhankelijk van of u zich in de {\em zinnebeeldmodus}\index{zinnebeeldmodus}\index{quote mode} bevindt.

De zinnebeeldmodus is slechts een fraaie benaming voor wat er gebeurt wanneer u een aanhalingsteken in de computer typt: niets, totdat u nog een aanhalingsteken typt of op \specialkey{RETURN} drukt. U zult zich herinneren dat we het eerder hadden over het probleem van vreemde symbolen die verschijnen wanneer u de cursor-toetsen gebruikt. Toen wilden we u niet afleiden, maar dat is een verschijnsel van de zinnebeeldmodus: in de zinnebeeldmodus tonen veel speciale toetsen een {\em zinnebeeld}\index{zinnebeeld} dat hen vertegenwoordigt, in plaats van hun normale werking te vervullen. Drukt u in de zinnebeeldmodus bijvoorbeeld de cursor-toets naar-links in, dan verschijnt het zinnebeeld \stw{Ƣ}. Drukt u de cursor-toets naar-rechts in, dan verschijnt een \stw{Ž}; bij omhoog een \stw{Ʊ}, bij omlaag een \stw{ű} en bij \specialkey{CLR\\HOME} een \stw{ų}. Houdt u \specialkey{SHIFT} ingedrukt en drukt u \specialkey{CLR\\HOME} in, dan verschijnt een \stw{Ƴ}.

Alternatieve termen:
\begin{itemize}
  \item De Nederlandstalige handleiding van de Commodore 64 uit 1982 spreekt over een codeteken in plaats van een zinnebeeld. De handleiding vermijdt het benoemen van de modus.
  \item In Engelstalige Mega65-documentatie wordt de zinnebeeldmodus de {\em quote mode} genoemd, naar het aanhalingsteken dat hem start.
\end{itemize}


\needspace{3cm}
Laten we dit dus gebruiken om de tweede regel het scherm te laten wissen wanneer hij de melding \stw{RAAD HET GETAL TUSSEN 1 EN 100} weergeeft. De eerste keer dat u het probeert, is het wat verwarrend, maar zodra u de slag te pakken heeft, is het vrij eenvoudig. Wat we uiteindelijk willen, is een regel die er als volgt uitziet:

\begin{screencode}
2 PRINT "ƳRAAD HET GETAL TUSSEN 1 EN 100"
\end{screencode}

Begin hiervoor met het intypen van \stw{2 PRINT "}. Houd vervolgens de \specialkey{SHIFT}-toets ingedrukt en druk op \specialkey{CLR\\HOME}. Uw regel zou er nu als volgt uit moeten zien: \stw{2 PRINT"Ƴ}. Is dat het geval, dan is het gelukt! U kunt de regel nu op de normale manier verder typen. Indien u dat gedaan heeft, kunt u zoals gebruikelijk het \stw{LIST}-commando gebruiken om na te gaan of u het programma met succes heeft gewijzigd. U zou uw gewijzigde regel met het \stw{Ƴ}-zinnebeeld moeten zien.

\screenshotwrap{dutch/images/getting-started/guess-number-2}

\needspace{2cm}
Voert u het programma nu uit door \stw{RUN} in te typen en zoals gebruikelijk op \specialkey{RETURN} te drukken, dan geeft de 2\textsuperscript{e} regel de computer de opdracht eerst het scherm te wissen en daarna de rest van de melding weer te geven, als volgt:

\screenshotwrap{dutch/images/getting-started/guess-number-3}

Dit verbergt de programmalijst voor de gebruiker, zodat hij niet onmiddellijk kan zien wat ons geheime getal is. We kunnen onze gok intypen, net als voorheen, maar net als voorheen keert de computer na één gok terug naar de \stw{READY.}-prompt. We zouden het prettig vinden als mensen meer dan één gok konden wagen, zonder dat zij hoeven te weten dat zij het programma opnieuw moeten uitvoeren.

\needspace{3cm}
Er zijn een paar manieren waarop we dit kunnen aanpakken. We hebben het \stw{FOR} --\stw{NEXT}-patroon al gezien. Met dat patroon kunnen we het programma de gebruiker een zeker aantal gokken laten geven. Als we het \stw{NEXT}-commando door een andere programmaregel zouden laten volgen, konden we de gebruiker zelfs vertellen wanneer hij te veel gokken heeft gewaagd. Laten we dus eens naar ons programma kijken en bedenken hoe we dat zouden kunnen doen. Hier is onze huidige programmalijst opnieuw:

\begin{screencode}
1 SN=23
2 PRINT "ƳRAAD EEN GETAL TUSSEN 1 EN 100"
3 INPUT "WAT IS UW GOK"; G
4 IF G<SN THEN PRINT "MIJN GETAL IS GROTER"
5 IF G>SN THEN PRINT "MIJN GETAL IS KLEINER"
6 IF G=SN THEN PRINT "GEFELICITEERD! U HEEFT MIJN GETAL GERADEN!"
\end{screencode}

Als we willen dat de gebruiker meerdere gokken kan wagen, moeten regels 2 tot en met 6 meerdere keren worden uitgevoerd. Dat maakt het er voor ons wat lastiger op, want we moeten een regel tussen regel 1 en 2 invoegen. Tenzij u echter een wiskundig genie bent, zijn er geen gehele getallen tussen 1 en 2 en de MEGA65 begrijpt regelnummers zoals 1.5 niet.

\needspace{4cm}
Gelukkig beschikt de MEGA65 over het
\stw{RENUMBER}\index{RENUMBER}\index{BASIC 65-commando's!RENUMBER}
\index{Regels!hernummeren}-commando. Dit commando kan alleen in de directe modus worden ingetypt. Bij uitvoering verandert het de regelnummers in het programma, zonder hun volgorde te wijzigen. De nieuwe nummers zijn gewoonlijk veelvouden van 10, zodat u tussenin veel vrije nummers overhoudt om extra regels toe te voegen. Als we het bijvoorbeeld op ons programma toepassen, dan hernummert het de regels tot 10, 20, \ldots, 60. Dat dit is gebeurd, kunnen we controleren met het \stw{LIST}-commando:

\screenshotwrap{dutch/images/getting-started/renumber-1}

\needspace{3cm}
Nu is ons leven een stuk eenvoudiger: we kunnen voor ons \stw{FOR}-commando elk getal tussen 10 en 20 kiezen. Het is een gebruikelijke keuze om het middelste getal te gebruiken, zodat u, als u later nog andere zaken wilt toevoegen, daarvoor de ruimte heeft. Laten we dus een \stw{FOR}-commando toevoegen dat de gebruiker 10 kansen geeft om het getal te raden. We kunnen hiervoor elke variabelenaam gebruiken die we willen, behalve \stw{G} en \stw{SN}, want die gebruiken we al. Het zou erg verwarrend zijn als we die door elkaar haalden! Laten we dus een regel als deze toevoegen:

\begin{screencode}
15 FOR I = 1 TO 10 STEP 1
\end{screencode}

\needspace{3cm}
Nu hebben we na regel 60 een bijpassend \stw{NEXT I} nodig. Laten we het handige patroon van het optellen van 10 aanhouden om het volgende regelnummer te bepalen en het als regel 70 toevoegen:

\begin{screencode}
70 NEXT I
\end{screencode}

\needspace{4cm}
Die regels kunnen we intypen, waarna we met het \stw{LIST}-commando controleren of het resultaat klopt:

\screenshotwrap{dutch/images/getting-started/guess-number-4}

Dat ziet er al heel aardig uit. Er zijn echter nog een paar kleine probleempjes. Kunt u bedenken welke? Wat gebeurt er nu nadat de gebruiker een gok waagt? En wat gebeurt er als zijn gokken op zijn?

Als u heeft bedacht dat het wagen van een gok het scherm onmiddellijk wist, dan mag u uzelf gerust op de schouder kloppen! De gebruiker heeft nauwelijks tijd om de melding te lezen. Erger nog: als hij het getal juist raadt, komt hij daar niet eens achter en het programma gaat gewoon door. We zouden graag zien dat het programma stopt of eindigt zodra de gebruiker een juiste gok maakt.

\index{STOP}\index{BASIC 65-commando's!STOP}
\index{END}\index{BASIC 65-commando's!END}
\index{CONT}\index{BASIC 65-commando's!CONT}
Dit kunnen we doen met het \stw{STOP}- of het \stw{END}-commando. Deze twee commando's lijken vrij veel op elkaar. Het belangrijkste verschil is dat de computer u vertelt waar het programma gestopt is wanneer u het met \stw{STOP} stopt en dat u de kans krijgt om het programma voort te zetten met het \stw{CONT}-commando. Het \stw{END}-commando vertelt de computer daarentegen dat het programma het einde bereikt heeft en dat de computer terug moet keren naar \stw{READY}. Voor ons programma is het \stw{END}-commando zinvoller, want nadat de gebruiker het getal geraden heeft, is er geen enkele reden meer om verder te gaan.

Nu hebben we een manier nodig om de computer twee verschillende dingen te laten doen wanneer de gebruiker een juiste gok maakt. We zouden gewoon een extra \stw{IF}-commando na regel 60 kunnen toevoegen dat de felicitatiemelding weergeeft, bijv. \stw{65 IF G=SN THEN END}.

We kunnen echter nog iets eleganter te werk gaan. Er is namelijk een manier om meerdere commando's op één enkele regel te zetten. Als u zich herinnert hoe we het \stw{INPUT}-commando hebben leren gebruiken, weet u misschien nog dat er twee verschillende tekens zijn waarmee stukjes gegevens van elkaar worden gescheiden: \stw{,} en \index{, (komma)}\stw{:}\index{: (dubbele punt)}. Het tweede teken, \stw{:}, heet een dubbele punt en kan ook gebruikt worden om BASIC-commando's op één regel van elkaar te scheiden. Willen we regel 60 dus zo wijzigen dat het \stw{PRINT}-commando de felicitatiemelding weergeeft en daarna het programma beëindigt met \stw{END}, dan kunnen we gewoon \stw{: END} aan het einde van de regel toevoegen. De regel zou er dan als volgt uit moeten zien:\index{: (dubbele punt)}

\begin{screencode}
  60 IF G=SN THEN PRINT "GEFELICITEERD! UW HEEFT MIJN GETAL GERADEN!": END
\end{screencode}

Dat probleem is daarmee opgelost. Het zou echter ook fijn zijn om het scherm niet telkens na een gok te wissen, zodat de gebruiker kan zien wat zijn laatste gok was en of die groter of kleiner was dan het getal. Daartoe kunnen we het wissen van het scherm uit regel 20 halen en een nieuw PRINT-commando aan een lager regelnummer toevoegen, zodat het scherm één keer aan het begin van het programma gewist wordt, voordat de gebruiker met gokken kan beginnen.

\needspace{2cm}
We zouden het bijvoorbeeld in regel 5 kunnen zetten, zodat het de allereerste handeling van het programma is. Zoals we eerder opmerkten, doen de regelnummers zelf er niet toe. Het enige dat telt, is te onthouden dat de computer bij het laagste regelnummer begint en de regels op volgorde uitvoert. Laten we die wijzigingen aan ons programma dan maar doorvoeren:

\begin{screencode}
  20 PRINT "RAAD EEN GETAL TUSSEN 1 EN 100"
  5 PRINT "J"
\end{screencode}

\needspace{4cm}
Als u die regels intypt en het programma opnieuw met \stw{LIST} weergeeft, zou u ongeveer het volgende moeten zien:

\screenshotwrap{dutch/images/getting-started/guess-number-5}

\needspace{4cm}
We kunnen het programma nu met \stw{RUN} uitvoeren en kijken of het werkt. Laten we het proberen!

\screenshotwrap{dutch/images/getting-started/guess-number-6}

Het scherm wordt nog steeds gewist en dat is goed. Ziet u echter al één klein verschil? Boven de eerste melding staat een lege regel. Dat komt doordat ons \stw{PRINT}-commando in regel 5 na het weergeven van het teken dat het scherm wist naar de volgende rij op het scherm gaat. Dat kunnen we verhelpen door een \stw{;} (puntkomma) aan het einde van het \stw{PRINT}-commando te plaatsen. Dat vertelt het \stw{PRINT}-commando dat het na alles te hebben weergegeven niet naar het begin van de volgende rij op het scherm moet gaan. Veranderen we regel 5 dus in \stw{5 PRINT "Ƴ";}, dan verdwijnt de lege ruimte boven aan het scherm.

\needspace{4cm}
Laten we terugkeren naar ons programma: we kunnen nu gokken en het programma vertelt ons of elke gok meer of minder is dan het juiste getal. Na tien gokken vraagt het programma niet meer om een nieuwe gok en keert het terug naar de \stw{READY.}-prompt, als volgt:

\screenshotwrap{dutch/images/getting-started/guess-number-7}

\needspace{2cm}
Het zou fijn zijn de gebruiker te vertellen wanneer zijn gokken op zijn. We moeten deze melding na het \stw{NEXT}-commando toevoegen. Het is ook netjes de gebruiker te vertellen wat het geheime getal was, in plaats van hem in onzekerheid te laten. Laten we de regel dus aan het einde van ons programma toevoegen, als regel 80:

\begin{screencode}
80 PRINT "HELAAS! UW KANSEN ZIJN OP. MIJN GETAL WAS"; SN
\end{screencode}

\needspace{4cm}
Raadt de gebruiker het getal nu niet, dan krijgt hij een nuttige melding, als volgt:

\screenshotwrap{dutch/images/getting-started/guess-number-8}

\subsection{Oefeningen}

{\bf 1. Kunt u het programma aan het begin naar het geheime getal laten vragen?}

Op dit moment zet het programma het geheime getal elke keer op 23. Om het spel interessanter te maken zou het mooi zijn de eerste gebruiker naar het geheime getal te vragen en daarna de rest van het spel te starten, zodat iemand anders kan proberen het getal te raden.

{\em Aanwijzing:} U zult de regel die de \stw{SN}-variabele instelt zo moeten wijzigen dat de waarde bij de eerste gebruiker kan worden opgevraagd. Het \stw{INPUT}-commando kan daarbij van pas komen.

{\bf 2. Kunt u het programma naar de naam van de gebruiker laten vragen en het vervolgens gepersonaliseerde antwoorden laten geven?}

Op dit moment toont het programma zeer eenvoudige meldingen. Het zou mooi zijn om de gebruiker naar zijn of haar naam te vragen en die naam vervolgens te gebruiken om gepersonaliseerde meldingen te maken, zoals \stw{HELAAS DIRK, MAAR DAT GETAL IS TE KLEIN}.

{\em Aanwijzing: U zult vroeg in het programma een regel moeten toevoegen om de gebruiker naar zijn of haar naam te vragen.}

{\em Aanwijzing: Mogelijk wilt u nog eens bekijken hoe we het \stw{PRINT}-commando hebben gebruikt, onder meer met \stw{;} om meer dan één waarde op een regel weer te geven.}

{\bf 3. Kunt u het uiterlijk van de meldingen verbeteren met kleuren en een betere spatiëring?}

Het programma is eigenlijk niet bepaald fraai geworden. Het zou mooi zijn om kleuren te gebruiken.

{\em Aanwijzing: Mogelijk wilt u meer \stw{PRINT}-commando's toevoegen om de spatiëring en de opmaak van de meldingen te verbeteren.}

{\em Aanwijzing: Mogelijk wilt u de kleurcodes gebruiken in de meldingen die u met \stw{PRINT} weergeeft.}

{\em Aanwijzing: Mogelijk wilt u ook de commando's \stw{FOREGROUND}, \stw{BACKGROUND} en \stw{BORDER} gebruiken om de kleur van de tekst, de schermachtergrond en de rand in te stellen.}
\index{FOREGROUND}\index{BASIC 65-commando's!FOREGROUND}
\index{BACKGROUND}\index{BASIC 65-commando's!BACKGROUND}
\index{BORDER}\index{BASIC 65-commando's!BORDER}


{\bf 4. Kunt u het programma laten aangeven of een gok ``warmer'' of ``kouder'' is dan de vorige gok?}

Op dit moment vertelt het programma u alleen of de gok hoger of lager is dan het geheime getal. Het zou mooi zijn als het programma u bij elke gok kon vertellen of de gok dichterbij komt of juist verder weg raakt: wanneer de gok dichterbij komt, moet het de gebruiker laten weten dat hij ``warmer'' wordt en ``kouder'' wanneer hij verder weg raakt.

Dit is heel wat ingewikkelder dan de voorgaande oefeningen en u zult een aantal nieuwe aangelenheden moeten uitdenken.

{\em Aanwijzing: U zult de vorige gok in een aparte variabele moeten onthouden en die vervolgens met de laatste vergelijken: is hij dichterbij of verder weg. Mogelijk heeft u \stw{IF}-commando's nodig waarin op het eerste \stw{IF} nog een ander volgt, of moet u leren de operator \stw{AND} te gebruiken.\index{AND}\index{BASIC 65-operatoren!AND}}

\section{Willekeurige getallen en toeval}

Op het spelletje ``raad het getal'' komen we zo terug, maar laten we eerst een uitstapje maken. Een uitstapje door een doolhof. Laten we hopen dat we er vóór het einde van de les weer uitkomen! Laten we een eenvoudige manier bekijken om een doolhof te maken. Dit programma is reeds lange tijd bekend. Het werkt door willekeurig te kiezen of het een {\symbolfont{M}}- of een {\symbolfont{N}}-teken weergeeft. Deze tekens krijgt u door \specialkey{SHIFT} ingedrukt te houden en op de toets N of M te tikken. U kunt de tekens op de voorkant van die toetsen zien. Hoewel ze op de toetsen binnen een kader staan afgebeeld, verschijnt dat kader niet; alleen de diagonale lijn is zichtbaar. Het blijkt dat het herhaald lukraak afdrukken van één van deze twee tekens een aardig doolhof oplevert.

\needspace{4cm}
Laten we het eens proberen. Om dit te kunnen doen, hebben we een manier nodig om willekeur te genereren. De MEGA65 heeft daarvoor de functie \stw{RND(1)}. Deze functie werkt als een variabele, maar telkens wanneer u haar gebruikt, levert zij een ander resultaat op. Laten we bekijken hoe dat in zijn werk gaat. Typ het volgende in:\index{RND()}\index{BASIC 65-functies!RND}

\begin{screencode}
PRINT RND(1)
\end{screencode}

Elke keer dat u dit intypt, levert het een ander antwoord op, zoals u hier kunt zien:

\screenshotwrap{images/getting-started/rnd-1}

We kunnen zien dat dit ons verschillende resultaten oplevert: \stw{1.07870447E-03},\stw{.793262171}, \stw{.44889513}, \stw{.697215893}. Elk van deze getallen ligt tussen 0 en 1, ook het eerste. Het eerste getal is in {\em wetenschappelijke notatie}\index{wetenschappelijke notatie} geschreven. De \stw{E-03} betekent dat de waarde $1.07870447\times10^{-3}$ = 0.000107870447 is. Dat wil zeggen dat de \stw{E-03} aangeeft dat de decimale punt drie plaatsen naar links verschuift. Staat er een \stw{+} achter de \stw{E}, dan betekent dat dat de decimale punt naar rechts verschuift. \stw{1.23456E+3} staat bijvoorbeeld voor het getal 1234.56.

\needspace{3cm}
Goed, we hebben u een doolhof beloofd, dus geven we u er maar één. Met deze \stw{RND(1)} kunnen we tussen deze twee tekens kiezen. Het eerste heeft een tekencode van 205 en het tweede heel handig 206. Dat betekent dat we, wanneer we het resultaat van \stw{RND(1)} bij 205.5 optellen, een getal tussen 205.5 en 206.5 krijgen. De helft van de tijd is dat $205.iets$ en de andere helft van de tijd $206.iets$. Zo kunnen we het ene of het andere teken weergeven met de functie \stw{CHR\$()}, welke het teken teruggeeft dat hoort bij het getal dat we tussen de haakjes zetten. Met die functie kunnen we iets doen als:

\begin{screencode}
  LET C = 205.5+RND(1)
  PRINT CHR$(C);
\end{screencode}

\needspace{3cm}
Dit geeft elke keer het ene of het andere teken weer. Hiermee konden we het doolhof al afdrukken door dit keer op keer te herhalen en er een lus van te maken. Daarvoor zouden we \stw{FOR} en \stw{NEXT} kunnen gebruiken. In dit geval willen we echter dat het programma voor altijd blijft doorgaan, dat wil zeggen: elke keer dat het programma aan het einde komt, moet het weer bij het begin beginnen. De makers van BASIC waren niet erg creatief, want het commando dat dit doet heet \stw{GOTO}\index{GOTO}\index{BASIC 65-commando's!GOTO}. U zet daarachter het nummer van de regel die vervolgens uitgevoerd moet worden, bijv. \stw{GOTO 1}. Hiermee kunnen we ons kleine doolhofprogramma zo schrijven dat het onafgebroken blijft draaien:

\begin{screencode}
10 LET C = 205.5+RND(1)
20 PRINT CHR$(C);
30 GOTO 10
\end{screencode}

\needspace{4cm}
Voert u dit programma uit met \stw{RUN}, dan begint de computer eindeloos een doolhof te tekenen, dat eruitziet zoals op de onderstaande schermafbeelding. U kunt het programma op elk moment stoppen door \specialkey{RUN\\STOP} in te drukken en u kunt het onderbreken door \specialkey{NO\\SCROLL} in te drukken en hervatten door opnieuw \specialkey{NO\\SCROLL} in te drukken\index{NO SCROLL}. Drukt u \specialkey{RUN\\STOP} in, dan vertelt de computer u waar hij op dat moment gebleven was. In het geval van de onderstaande schermafbeelding was hij bezig met regel 10:

\screenshotwrap{images/getting-started/rnd-2}

\needspace{2cm}
Dat werkt prima en tekent een zeer beroemd doolhof \cite{montfort201210}. Toch kunnen we het programma kleiner maken. Het resultaat van de berekening van welk symbool er weergegeven moet worden, hoeft niet op een aparte regel te staan. We kunnen de berekening rechtstreeks tussen de haakjes van de functie \stw{CHR\$()} zetten:

\begin{screencode}
10 PRINT CHR$(205.5+RND(1));
20 GOTO 10
\end{screencode}

\needspace{2cm}
En we kunnen gebruikmaken van wat we over het \stw{:}-symbool (de dubbele punt) hebben geleerd en het \stw{GOTO}-commando op dezelfde regel zetten als het \stw{PRINT}-commando:

\begin{screencode}
10 PRINT CHR$(205.5+RND(1));: GOTO 10
\end{screencode}

Ziet u hoe er vaak vele manieren zijn om met een programma hetzelfde effect te bereiken? Dat is heel normaal. Bij complexe programma's zijn er heel veel manieren om hetzelfde effect te bereiken. Dit is één van de terreinen in het computerprogrammeren waarop u zeer creatief kunt zijn.

Laten we terugkeren naar het onderwerp willekeur. Het is goed dat we deze willekeurige getallen tussen 0 en 1 kunnen gebruiken om een doolhof te tekenen, maar het is wat lastig om mensen te vragen een heel lang decimaal getal in te typen. Willen we een getal in het bereik 1 tot 100, dan kunnen we wat we van \stw{RND(1)} krijgen met 100 vermenigvuldigen. Doen we dat, dan wordt het al wat beter, maar we krijgen nog steeds getallen als \stw{55.0304651}, \stw{30.3140154}, \stw{60.2505497} en \stw{.759229916}.

\needspace{3cm}
Dat komt al in de buurt, maar we willen die gebroken delen echt kwijt. We willen namelijk hele getallen, ofwel {\em integers}. BASIC heeft daarvoor de functie \stw{INT()}, welke net zo werkt als de functie \stw{RND(1)}: welk getal u ook tussen de haakjes zet, de functie geeft alleen het gehele deel daarvan terug. \stw{INT(2.18787)} geeft bijvoorbeeld de waarde \stw{2} terug. Zoals we zojuist zeiden, hakt de functie het gebroken deel eraf, dat wil zeggen dat hij altijd naar beneden afrondt. Zelfs \stw{INT(2.9999999)} levert dus nog steeds \stw{2} op in plaats van \stw{3}. Dat betekent dat we, als we het resultaat van \stw{RND(1)} met 100 vermenigvuldigen, een getal in het bereik 0 -- 99 krijgen in plaats van 1 -- 100. Dat is gelukkig eenvoudig op te lossen: we tellen gewoon 1 bij het resultaat op. Om dus een integer te genereren, dat wil zeggen een geheel getal tussen 1 en 100, beide inbegrepen, kunnen we iets doen als:

\begin{screencode}
  PRINT INT(RND(1)*100) + 1
\end{screencode}

\needspace{3cm}
Dat ziet er veel beter uit. Laten we ons ``raad het getal''-programma nogmaals intypen. Deze keer vervangen we echter de plek waar we ons geheime getal op 23 zetten door een willekeurig geheel getal tussen 1 en 100. Kijk nog niet naar de oplossing. Denk eerst na over hoe u met het bovenstaande \stw{SN} op een willekeurig geheel getal tussen 1 en 100 kunt zetten. Heeft u uw eigen gok klaar, bekijk dan wat wij hieronder hebben bedacht. Wellicht heeft u een ander programma gemaakt dat hetzelfde kan doen. Dat is ook prima!

\begin{screencode}
10 SN=INT(RND(1)*100)+1
20 PRINT "Ƴ"
30 FOR I = 1 TO 10 STEP 1
40 PRINT "RAAD HET GETAL TUSSEN 1 EN 100"
50 INPUT "WAT IS UW GOK"; G
60 IF G<SN THEN PRINT "MIJN GETAL IS GROTER"
70 IF G>SN THEN PRINT "MIJN GETAL IS KLEINER"
80 IF G=SN THEN PRINT "GEFELICITEERD! U HEEFT MIJN GETAL GERADEN!": END
90 NEXT I
100 PRINT "HELAAS, UW KANSEN ZIJN OP"
\end{screencode}

Nu hoeven we niet langer te vrezen dat iemand het getal raadt en we hoeven ook niemand anders het getal voor ons te laten kiezen. Daardoor is het programma veel leuker om te spelen. Kunt u het verslaan?


\subsection{Oefeningen om te proberen}

{\bf 1. Kunt u het doolhofprogramma verschillende doolhoven laten tekenen?}

Het doolhofprogramma toont momenteel evenveel {\symbolfont{N}} als {\symbolfont{M}}. Kunt u het programma zo veranderen dat het twee keer zoveel van het ene teken afdrukt als van het andere? Hoe ziet het doolhof er dan uit?

{\em Aanwijzing: Wij hebben \stw{205.5} gebruikt, zodat we na het optellen van een willekeurig getal tussen 0 en 1 de helft van de tijd op $205.iets$ en de andere helft van de tijd op $206.iets$ uitkomen. Verlaagt u \stw{205.5} richting \stw{205}, of verhoogt u het richting \stw{206}, dan verandert u de relatieve verhouding waarin elk teken verschijnt.}

{\bf 2. Kunt u het ``raad het getal''-programma aanpassen zodat het een getal tussen 1 en 10 kiest?}

Op dit moment kiest het programma een getal tussen 1 en 100. Pas het programma aan zodat het een getal uit een ander bereik kiest. Vergeet niet de melding aan te passen die de computer aan de gebruiker toont. Heeft hij dan nog steeds 10 gokken nodig? Verander het maximum aantal gokken dat hij krijgt voordat hij verliest in een geschikter aantal.

{\em Aanwijzing: U zult de regel moeten aanpassen die \stw{SN} instelt, evenals de \stw{PRINT}-melding die de gebruiker uitleg geeft.}

{\bf 3. Zet de scherm-, rand- en tekstkleur op willekeurige kleuren}

Pas het doolhofprogramma of het ``raad het getal''-programma aan om willekeurige kleuren te gebruiken. Hoe kunt u ervoor zorgen dat de tekst altijd zichtbaar blijft?

{\em Aanwijzing: Gebruik de commando's \stw{FOREGROUND}, \stw{BACKGROUND} en \stw{BORDER} om de kleuren in te stellen. Gebruik kleurnummers tussen 0 en 15, inclusief. U kunt aan het einde van deze commando's een berekening plaatsen in plaats van een eenvoudig getal.}

{\em Aanwijzing: Om te voorkomen dat u de tekstkleur op dezelfde kleur als de achtergrond instelt, kunt u berekenen welke achtergrondkleur u wilt gebruiken en die in de ene variabele bewaren. Bereken vervolgens de tekstkleur die u wilt gebruiken en bewaar die in een andere variabele. Hebben de twee variabelen hetzelfde getal, dan moet u een van beide veranderen.}

{\bf 4. Laat het ``raad het getal''-programma bij het starten willekeurig kiezen tussen twee verschillende welkomstmeldingen}

Het ``raad het getal''-programma toont momenteel telkens wanneer het start dezelfde melding. Pas het aan zodat het elke keer een van twee mogelijke meldingen toont.

{\em Aanwijzing: Gebruik \stw{RND(1)} om een willekeurig getal te verkrijgen. Is dat getal kleiner dan een bepaalde drempelwaarde, toon dan de eerste melding; zo niet, toon dan de tweede melding.}

{\em Aanwijzing: Het kan gemakkelijker zijn om het willekeurige getal in een variabele te bewaren, zodat u twee \stw{IF}-instructies kunt gebruiken om te bepalen of u elke melding toont.}

{\em Aanwijzing: Gebruikt u \stw{<} (kleiner dan) als relationele operator in een van de \stw{IF}-instructies, dan moet u in de andere de tegenovergestelde gebruiken. Het tegenovergestelde van kleiner dan is groter dan of gelijk aan.}
